\documentclass[11pt]{article}
\usepackage{amsmath,amssymb,amsthm}
\usepackage{mathtools}

\theoremstyle{definition}
\newtheorem{definition}{Definition}
\newtheorem{assumption}{Assumption}
\newtheorem{property}{Empirical Property}
\newtheorem{remark}{Remark}
\newtheorem{proposition}{Design Implication}

\newcommand{\Route}{\mathrm{R}}
\newcommand{\Place}{\mathrm{P}}
\newcommand{\Traffic}{\mathrm{T}}

\title{Formalizing MoE Routing, Placement, and Induced Traffic}
\author{}
\date{}

\begin{document}
\maketitle

This note formalizes three claims used to argue against treating MoE-induced
network traffic as unstructured: (i) routing is a deterministic function of
input and model, not a stochastic process; (ii) routing is, by construction,
independent of expert placement; and (iii) the traffic that placement induces
is nonetheless empirically local, which is the distinct fact that licenses
regionally-reconfigurable optical interconnects.

\section{Setup}

\begin{definition}[Model and input space]
Let $\theta$ denote the frozen parameters of a trained MoE model, $\mathcal{X}$
the space of input sequences, and $E=\{e_1,\dots,e_K\}$ the set of $K$ experts
at a given MoE layer. For $x\in\mathcal{X}$, let $h_l(x_t\mid\theta)\in\mathbb{R}^d$
denote the hidden representation of token $x_t$ at layer $l$, computed by the
deterministic preceding sub-network of the model.
\end{definition}

\begin{definition}[Router]
The router at layer $l$ is
\[
  \Route_l:\mathbb{R}^d\times\Theta_l\longrightarrow\Delta^{K-1},\qquad
  \Route_l\bigl(h_l(x_t\mid\theta);\theta_l\bigr)=\mathrm{softmax}\bigl(W_l\,h_l(x_t\mid\theta)+b_l\bigr),
\]
followed by a $\mathrm{top}\text{-}k$ operator mapping $\Delta^{K-1}$ to a subset
$S_{l,t}(x)\subseteq E$, $|S_{l,t}(x)|=k$.
\end{definition}

\section{Determinism of routing}

\begin{assumption}[Routing as a deterministic function of input and model]
\label{ass:determinism}
For fixed, frozen $\theta$ and given $x\in\mathcal{X}$, the routing decision at
every layer and token position is a deterministic function of $(x,\theta)$ alone:
\[
  S_{l,t}(x) = F_l(x,\theta), \qquad\text{i.e.}\qquad
  \mathrm{Routing} = F(\mathrm{inputs},\ \mathrm{model}).
\]
Routing is \emph{not} an i.i.d.\ stochastic process over $(l,t)$. Any observed
run-to-run variability is attributable to factors external to $F$ itself
(batch-dependent capacity dropping, non-associative floating-point reduction
under distributed all-reduce, or training-time-only noise injection).
\end{assumption}

\begin{remark}[Latent conditioning exposes task/domain structure]
\[
  S_{l,t}(x) = F_l\bigl(h_l(x_t\mid z(x)),\,\theta_l\bigr), \qquad z(x)\in\mathcal{Z},
\]
where $z(x)$ is a latent task/domain/language variable. Empirically,
$\mathrm{Corr}(S_{l,t},S_{l,t-1})$ and $\mathrm{Corr}(S_{l,t},S_{l-1,t})$ are
bounded away from $0$, and $P(S_{l,t}\mid z)$ is markedly more concentrated
than the pooled marginal $P(S_{l,t})$.
\end{remark}

\section{Independence from placement}

\begin{definition}[Placement]
A placement is an assignment $\pi:E\to\mathcal{D}$ mapping experts to physical
devices $\mathcal{D}=\{d_1,\dots,d_M\}$, chosen by the parallelization/sharding
strategy and set independently of $\theta$.
\end{definition}

\begin{assumption}[Routing--placement independence]
\label{ass:independence}
For a communication-oblivious router (Assumption~\ref{ass:determinism} with
$F_l$ having no dependence on $\pi$),
\[
  \forall\,\pi,\pi'\in\Pi:\qquad
  P\bigl(S_{l,t}(x)\mid\theta,\pi\bigr) = P\bigl(S_{l,t}(x)\mid\theta,\pi'\bigr).
\]
The logical routing distribution is invariant under any relabeling of
expert-to-device placement.
\end{assumption}

\begin{definition}[Logical demand matrix]
Fix a layer $l$. Let $A_l(x,\theta)\in\mathbb{N}^{|\mathcal{D}|\times K}$ be the
matrix whose $(a,e)$ entry counts token-expert selections originating at
device $d_a$ and targeting expert $e$:
\[
  A_l(x,\theta)_{a,e} \;=\; \bigl|\{\,t : \mathrm{src}(x_t)=d_a,\ e\in S_{l,t}(x)\,\}\bigr|.
\]
Here $\mathrm{src}(x_t)\in\mathcal{D}$ is the device currently holding token
$x_t$, fixed by the (non-expert) tensor/data-parallel sharding and treated as
exogenous to this section. By Assumption~\ref{ass:determinism},
$A_l(x,\theta)$ depends only on $(x,\theta)$ — never on $\pi$.
\end{definition}

\begin{definition}[Induced physical traffic]
\label{def:traffic}
Let $\Pi_\pi\in\{0,1\}^{K\times|\mathcal{D}|}$ be the placement matrix with
$(\Pi_\pi)_{e,d}=1$ iff $\pi(e)=d$ (one $1$ per row: each expert lives on one
device). The induced device-to-device traffic is simply the matrix product
\[
  \Traffic_\pi(x,\theta) \;=\; A_l(x,\theta)\,\Pi_\pi \;\in\;\mathbb{N}^{|\mathcal{D}|\times|\mathcal{D}|},
  \qquad
  \Traffic_\pi(d_a,d_b) = \sum_{e} A_l(x,\theta)_{a,e}\,(\Pi_\pi)_{e,b}.
\]
$A_l$ carries all the input/model dependence; $\Pi_\pi$ carries all the
placement dependence; the traffic matrix is their product.
\end{definition}

% \begin{definition}[Induced physical traffic]
% Given a routing trace $\{S_{l,t}(x)\}$ and placement $\pi$, the induced traffic
% between devices $d_a,d_b\in\mathcal{D}$ is
% \[
%   \Traffic_\pi(d_a\to d_b) = \sum_{l,t,x,e}
%   \mathbb{1}[\mathrm{src}(x_t)=d_a]\cdot\mathbb{1}[e\in S_{l,t}(x)]\cdot\mathbb{1}[\pi(e)=d_b],
% \]
% so that
% \[
%   \Traffic = \Route(x,\theta)\ \circ\ \Place(\pi)
% \]
% factorizes into a placement-independent logical term and a routing-independent
% physical term.
% \end{definition}

\section{Locality: empirical, not definitional}

\begin{property}[Bounded-locality of induced traffic]
\label{prop:locality}
Let $B(e)$ denote the MoE block / expert-parallel group containing expert $e$,
and $\mathcal{D}_{B(e)}\subset\mathcal{D}$ the device subset hosting that group,
with $|\mathcal{D}_{B(e)}|\ll|\mathcal{D}|$. Empirically, for a fixed
parallelization strategy,
\[
  \Pr\bigl[\pi(e)\in\mathcal{D}_{B(e)}\bigr]\approx 1,
\]
i.e., although $S_{l,t}$ has full support over $E$ in principle, $\Traffic_\pi$
concentrates within bounded device subsets rather than spreading uniformly
over $\mathcal{D}\times\mathcal{D}$.
\end{property}

\begin{remark}
Assumption~\ref{ass:independence} and Property~\ref{prop:locality} are
logically independent. The former says the \emph{router} does not consult
$\pi$; the latter says the \emph{resulting} $\Traffic_\pi$ happens to be
spatially concentrated. Both are needed to justify regional, rather than
global, reconfigurability of the interconnect.
\end{remark}

\section{Design implication}

\begin{proposition}[Regional reconfigurability suffices]
Under Assumptions~\ref{ass:determinism}--\ref{ass:independence} and
Property~\ref{prop:locality}, the interconnect fabric $\mathcal{F}$ need only
support dynamic reconfiguration within each $\mathcal{D}_{B(e)}$, not across
all of $\mathcal{D}\times\mathcal{D}$:
\[
  \mathcal{F} = \mathcal{F}_{\mathrm{electrical}}\bigl(\text{deterministic }TP/DP/PP\bigr)
  \ \cup\ \mathcal{F}_{\mathrm{optical}}\bigl(\{\mathcal{D}_{B(e)}\}_{e\in E}\bigr).
\]
\end{proposition}

\section{Scope boundary}

\begin{remark}[Communication-aware routing breaks Assumption~\ref{ass:independence}]
If the router is trained with a locality-aware auxiliary loss or device-limited
capacity constraint, $F_l$ is no longer independent of $\pi$:
\[
  S_{l,t}(x) = F_l(x,\theta,\pi),
\]
and the factorization $\Traffic=\Route\circ\Place$ no longer holds; placement
and routing must then be co-optimized rather than solved in two stages.
\end{remark}

\end{document}